\documentclass[11pt]{article}
\usepackage{preamble}
\renewcommand{\answer}[1]{}

\begin{document}

\ladderheading{AP Statistics \textperiodcentered\ Topic 6.4 (CED 3.5) \textperiodcentered\ B: 2026-12-03 \textperiodcentered\ E: 2027-01-08}{Which Test Answers the Question Asked?}

\focusbanner{TODAY'S PROBLEM}

A district of 2{,}400 students had bus ridership $p_0=0.40$ before a new route. Before collecting data, the director asks whether the current proportion $p$ has \emph{increased}. An SRS of 180 current students finds 84 riders.\par\smallskip \textbf{Maya:} ``Use $H_0:p=0.40$ versus $H_a:p>0.40$ and a one-sample $z$-test. Check large counts with 72 and 108.''\par \textbf{Theo:} ``Use $H_a:p\neq0.40$ because samples can differ either way. Check large counts with the observed 84 and 96.''\par
\begin{center}
\textbf{Current SRS counts}\par\smallskip
\begin{tabular}{|l|c|c|c|}
\hline
\textbf{Group} & \textbf{Bus} & \textbf{Other} & \textbf{Total} \\ \hline
\textbf{Students} & 84 & 96 & 180 \\ \hline
\end{tabular}
\end{center}
\begin{firsttakebox}{FIRST TAKE --- before the video (5 min)}
Which setup better answers whether ridership increased? Choose Maya or Theo and cite one detail from the question or sample.
\workspace{0.9in}
\end{firsttakebox}

\begin{callout}{\IconBook\ SET UP THE QUESTION BEFORE USING THE SAMPLE (keep on the page)}{calloutgreen}
Define $p$ as the \textbf{population proportion}. For a benchmark $p_0$, use $H_0:p=p_0$.
Choose $H_a:p>p_0$, $p<p_0$, or $p\neq p_0$ from the research question stated \emph{before} seeing data;
$\hat p$ is sample evidence and never belongs in the hypotheses. For one categorical variable and one sample,
use a one-sample $z$-test for a population proportion. Verify random data, $n\leq0.10N$ when sampling without
replacement, and the null-expected counts $np_0\geq10$ and $n(1-p_0)\geq10$.
\end{callout}

\newpage
\textbf{Finish after the video:}\par\smallskip

\dokbadge{1}~\textbf{(a)} Define the population parameter $p$ in context and identify the null value $p_0$.\par
\workspace{0.7in}

\dokbadge{2}~\textbf{(b)} Compute $\hat p$, name the testing procedure, and verify the random, 10 percent, and large-counts conditions using the appropriate values.\par
\workspace{1.4in}

\dokbadge{3}~\textbf{(c)} Choose the warranted setup. Use the word \emph{increased}, distinguish $p_0$ from $\hat p$, and explain why large counts uses null-expected values. Give the reader consequence of Theo's alternative and count check.\par
\workspace{2.0in}

\begin{sentenceframebox}\raggedright\textbf{Frame:} The warranted alternative is $H_a:p\blank[0.4in]0.40$ because the director asked whether \blank[2.2in].\end{sentenceframebox}

\begin{sentenceframebox}\raggedright\textbf{Frame:} The large-counts check uses \blank[0.8in] and \blank[0.8in], not \blank[1.0in], because \blank[2.4in].\end{sentenceframebox}

\begin{summaryexitbox}{EXIT --- turn this in}
One thing the video changed about my first take: \hrulefill\par
\workspace{0.4in}
\end{summaryexitbox}

\end{document}
